\section{NN-S005: special values of every shifted Dedekind sheet}
\label{sec:NN-S005}

\paragraph{Inputs.} Corpus 01, Sections 5--8, and corpus 02,
Section 14, give the shift mechanism. The residue calculation below
extends the displayed Hurwitz special-value formula to every number
field and every real shift $t\geq0$. It uses the meromorphic
continuation of the Dedekind zeta function, with its sole simple
pole at one. It uses neither A nor M.

\begin{theorem}[NN-S005]
For a number field $K$, set $\kappa_K=\operatorname{Res}_{s=1}\zeta_K(s)$
and
\[
 F_{K,t}(s)=\sum_{0\ne J\subseteq\mathcal O_K}(N(J)+t)^{-s},
 \qquad t\geq0,\quad \Re s>1.
\]
This function continues meromorphically to $\mathbb C$, with only a
simple pole at one and residue $\kappa_K$. For every integer $m\geq0$,
\[
 \boxed{F_{K,t}(-m)=
 \sum_{k=0}^m\binom mk t^k\zeta_K(-m+k)
       -\frac{\kappa_K}{m+1}t^{m+1}.}                 \tag{S005.1}
\]
In particular $F_{K,t}(0)=\zeta_K(0)-\kappa_Kt$.
The last term is essential: substituting $s=-m$ before resolving
the zero times pole in the shift expansion would omit it.
\end{theorem}

\begin{proof}
Write $\zeta_K(s)=\sum_{n\geq1}a_n n^{-s}$ and choose an integer
$M>t$. Separate the finitely many terms with $n\leq M$. The binomial
series, absolutely convergent for $t/n<1$, gives
\[
 F_{K,t}(s)=\sum_{n\leq M}a_n(n+t)^{-s}
  +\sum_{k\geq0}\frac{(-t)^k(s)_k}{k!}
       \left(\zeta_K(s+k)-\sum_{n\leq M}a_n n^{-s-k}\right).       \tag{S005.2}
\]
First this equality holds in $\Re s>1$ by absolute convergence.
For a compact set of $s$-values, $(s)_k/k!$ is bounded by a fixed
power of $k$ times a constant. Choose a fixed $c$ large enough that
$\sum a_n n^{-\Re s-c}$ converges uniformly on that compact set.
For $k\geq c$ the parenthesis in (S005.2), expressed by its
convergent tail series, is bounded by a constant times
$(M+1)^{-k+c}$. Thus the $k$-series converges normally away from
its initially displayed poles. The term indexed by $k\geq1$ has
an apparent pole at $s=1-k$, which is cancelled by its simple
factor $s+k-1$. Only the $k=0$ pole at one remains, with the stated
residue. This also justifies evaluation of the continued sum.

At $s=-m$, terms with $k\leq m$ have coefficient
$(-t)^k(-m)_k/k!=\binom mk t^k$. The term $k=m+1$ has limit
\[
 \frac{(-t)^{m+1}}{(m+1)!}
 \left(\prod_{j=0}^{m-1}(-m+j)\right)\kappa_K
       =-\frac{\kappa_K t^{m+1}}{m+1}.
\]
Every term $k>m+1$ vanishes because the rising factorial is zero
and the zeta factor is regular. The finite subtractions in
(S005.2) cancel its initial finite sum by
$(n+t)^m=\sum_{k=0}^m\binom mk t^k n^{m-k}$.
This proves (S005.1), including $m=0$ with the empty product equal
to one.
\end{proof}

\paragraph{Applications and precise limits.}
NN-S004 identifies the separated sheet for a support lattice of
size $\ell$ with $t=\ell-1$, so (S005.1) applies for arbitrary
finite support size. For $K=\mathbb Q$ and $t=1$ it agrees with
$\zeta(s)-1$; this particular identity does not extend to a
general $K$. For example $\zeta_{\mathbb Q(i)}(s)=\zeta(s)\beta(s)$,
$\beta(0)=1/2$, and $\kappa_{\mathbb Q(i)}=\pi/4$ give
$F_{\mathbb Q(i),1}(0)=-(1+\pi)/4$. These values do not assert
the location of the zeros of the shifted sheet. Neither A nor M
provides its missing functional equation.

\paragraph{Geometric dictionary for NN-S005.}
Each finite amplitude quotient is the arithmetic closed
subscheme $\operatorname{Spec}(\mathcal O_K/J)$; the lattice
changes its chosen quotient-cardinality statistic from $N(J)$ to
$N(J)+\ell-1$. The theorem is an analytic statement about that
statistic. It does not identify it with an Euler product or with
a trace of a finite Weil object. In the ordinary scheme dictionary
of EGA I, Section 4 (closed subschemes), $J$ remains the same
ideal sheaf; adjoining support classes is extra semiring data.
The residue term is analytic boundary data, not an additional
effective summand in $K_0$.

\section{NN-S006: all local phase atoms forced by reflection}
\label{sec:NN-S006}

\paragraph{Inputs and attribution.}
We use A to put the nontrivial zeros on the critical line and
corpus 10, Theorem \texttt{thm:local-phase-packet} and Proposition
\texttt{prop:packet-symmetry-local}, for the reduced logarithmic
germ. The phase formalism retains the source's attribution:
J.~Tobley, \emph{Unpublished correspondence on orthogonal-walk
phase decompositions, Taylor-series atomization, and affine
cyclotomic charts}, ``project source transcripts communicated with
the manuscript materials, May 2026.'' The following
all-orders restriction is the consequence of the supplied theorem
A; it assumes no simplicity.

\begin{theorem}[NN-S006]
Let $H$ be holomorphic near a point $\rho$ on $\Re s=1/2$, and
suppose $H(s)=\overline{H(1-\bar s)}$ there. If its zero at $\rho$
has order $m$, write
\[
 H(\rho+z)=\lambda z^m\exp h(z),\qquad
 h(0)=0,\quad h(z)=\sum_{n\geq1}c_nz^n.
\]
Then
\[
 \lambda=(-1)^m\bar\lambda,\qquad
 c_n=(-1)^n\bar c_n.                                  \tag{S006.1}
\]
In particular all tangent phase atoms
$\operatorname{Im}(c_ni^n)$ vanish. The normal logarithmic modulus
of the reduced germ is even in the real displacement, and its
normal phase, using the logarithm with value zero at the origin,
is odd.
This applies to every zero of $\xi$ under A. For a primitive
Dirichlet character with root number $\varepsilon$, it applies
to $H_\chi=\varepsilon^{-1/2}\Lambda(s,\chi)$ at every nontrivial
zero under A. For a Dedekind zeta completion it applies under A
when $K/\mathbb Q$ is abelian, and under M, conditionally, for
general $K$.
\end{theorem}

\begin{proof}
Since $1-\bar\rho=\rho$, reflection says
$H(\rho+z)=\overline{H(\rho-\bar z)}$. Comparing the leading
terms gives the first identity. Dividing by the leading terms
gives $\psi(z)=\overline{\psi(-\bar z)}$ for the reduced germ
$\psi=H(\rho+z)/(\lambda z^m)$ with $\psi(0)=1$.
The two holomorphic logarithms with value zero at the origin
are therefore equal: $h(z)=\overline{h(-\bar z)}$.
Coefficient comparison proves the second identity.
Thus even coefficients are real and odd coefficients purely
imaginary. Multiplication by $i^n$ makes either type real, giving
the tangent assertion; evaluating at real $z$ gives the other two.

For $\xi$ the reflection follows from its functional equation
and real coefficients. For primitive $\chi$ the two standard
identities are
$\Lambda(s,\chi)=\varepsilon\Lambda(1-s,\bar\chi)$ and
$\overline{\Lambda(\bar s,\chi)}=\Lambda(s,\bar\chi)$.
Choose a unit square root $\alpha=\varepsilon^{1/2}$.
Since $\bar\alpha=\alpha^{-1}$, the normalized function
$\alpha^{-1}\Lambda(s,\chi)$ has the stated reflection.
The Dedekind completion is self-dual with real coefficients.
The appropriate RH input is used only to put every zero in the
fixed line; the local argument itself applies to any zero already
known to lie there. It does not assert that a zero is simple or
that an individual deformation velocity is nonzero.
\end{proof}

\paragraph{Geometric and GCT dictionary for NN-S006.}
At each finite jet order $r$, the completed local coordinate
algebra is $\mathbb C[z]/(z^{r+1})$ with the semilinear involution
$z\mapsto-z$, conjugating scalars. With $w=-iz$ its invariant
real form is $\mathbb R[w]/(w^{r+1})$, and scalar extension to
$\mathbb C$ recovers the original algebra. Identity (S006.1)
is exactly descent of the reduced logarithmic section to this
real Artin scheme. The proof is coefficientwise descent, so no
unproved effectivity assertion is needed. The infinitesimal
neighbourhood convention is EGA I, Section 4, and finite locally
free descent is the usual affine faithfully flat descent
dictionary. This is a real form and a reflection symmetry;
it supplies no finite-field Frobenius operator and no positive
canonical basis for a nonstandard quantum group.

\section{NN-S007: the literal zero divisor and tail are not finite Weil objects}
\label{sec:NN-S007}

\begin{theorem}[NN-S007]
The literal infinite zero-support and tail constructions in
corpus 05--10 cannot themselves be finite constructible
algebraic Weil objects by the following two direct routes.
First, the set of nontrivial complex zeta zeros is not a
Zariski-constructible subset of $\mathbb A^1_{\mathbb C}$.
Second, if $V_n$ are nonzero vector spaces over a characteristic-zero
field for infinitely many $n$, then
\[
 \prod_nV_n\big/\bigoplus_nV_n
\]
has infinite dimension. A puts the zero support on a line but
removes neither obstruction.
\end{theorem}

\begin{proof}
A closed proper algebraic subset of the affine line over
$\mathbb C$ is finite, since a nonzero one-variable polynomial
has finitely many roots. A constructible subset is a finite union
of locally closed subsets, hence is finite or cofinite. The
nontrivial zero set is infinite, while its complement contains
every real $s>1$ and is infinite; it is therefore not constructible.

Choose indices $i_1<i_2<\cdots$ with $V_{i_n}\ne0$ and
choose $v_{i_n}\ne0$. For distinct positive integers
$a_1,a_2,\ldots$, put $a_j^nv_{i_n}$ in factor $i_n$ and zero
in every other factor, and consider the resulting classes. If a finite linear
combination is zero in the quotient, then
$\sum_{j=1}^rc_ja_j^n=0$ for all sufficiently large $n$.
Taking $r$ consecutive such $n$ gives an invertible Vandermonde
matrix, multiplied by a nonsingular diagonal matrix. Hence all
$c_j=0$. Infinitely many independent classes result. The same
argument applies if the original index set has further factors.
\end{proof}

\paragraph{Exact scope of the obstruction.}
Corpus 06, Theorem 6.8, and corpus 07's vector-tail proposition
identify the relevant tail as this product modulo finite support.
There is no objection to the analytic coherent sheaf
$\mathcal O_{\mathbb C}/(\xi)$ on the analytic affine line.
The assertion is about its proposed direct identification with a
finite constructible algebraic object, or with one finite-dimensional
Weil representation. A different finite realization of selected
data is not ruled out; it must be constructed and shown to retain
the desired structure constants. That is the precise surviving
gate in prompt 10(b).

\paragraph{Geometric dictionary for NN-S007.}
Constructibility requires a finite algebraic stratification and
finite-rank stalks in the SGA 4 constructible category. The first
proof excludes the proposed literal support; the second excludes
the proposed literal stalk at infinity. Ordinary finite-dimensional
Weil representations and the weight definition of Weil II 1.2.2
cannot be assigned to that infinite tail without an additional
finite quotient or realization. The Boolean skeleton's scalar
linearization instead vanishes by NN-S003. These are distinct
operations, neither of which creates the missing finite object.
